\documentclass[11pt]{amsart}
\usepackage[margin=1in]{geometry}
\usepackage{amsmath,amssymb,amsthm,mathtools}
\usepackage[T1]{fontenc}
\usepackage{lmodern}
\usepackage{microtype}
\usepackage{enumitem}
\usepackage[hidelinks]{hyperref}

\newtheorem{theorem}{Theorem}[section]
\newtheorem{lemma}[theorem]{Lemma}
\newtheorem{proposition}[theorem]{Proposition}
\newtheorem{corollary}[theorem]{Corollary}
\theoremstyle{definition}
\newtheorem{definition}[theorem]{Definition}
\newtheorem{example}[theorem]{Example}
\theoremstyle{remark}
\newtheorem{remark}[theorem]{Remark}

\newcommand{\R}{\mathbb{R}}
\newcommand{\Z}{\mathbb{Z}}

\title[Topological Capacity Veto in $D{=}3$]{Alexander-Duality Linking, Link Penalties,\\
and the Finite-Capacity Veto in Three Dimensions}

\author{Jonathan Washburn}
\address{Austin, Texas, USA}
\email{washburn.jonathan@gmail.com}
\date{February 2026}

\begin{document}
\begin{abstract}
We prove a conditional finite-capacity veto for rigid-body rotation as a
candidate blow-up profile of finite-energy ($H^1$) divergence-free data in
three dimensions.
The argument combines three ingredients:
(i)~\emph{Alexander duality}, which shows that an integer-valued linking
number for disjoint oriented loops exists if and only if the ambient
dimension is~$3$---this is a purely topological fact, independent of PDE
considerations;
(ii)~a \emph{link penalty}: each topological crossing of vortex
lines incurs a strictly positive cost $\Delta J=\ln\varphi>0$
(where $\varphi=(1+\sqrt{5})/2$ is the golden ratio), derived from the
canonical reciprocal cost~\cite{F1};
(iii)~a \emph{finite-capacity veto}: the $H^1$ energy of the initial data
provides a finite budget, which cannot fund the infinitely many link
crossings required under an explicit crossing-count hypothesis.
The main theorem (Theorem~\ref{thm:veto}) is stated with explicit
assumptions, clearly separating the unconditional topological and
analytic ingredients from the crossing-count input that must be verified
externally.
\end{abstract}

\subjclass[2020]{Primary 57K10; Secondary 76D05, 55M05, 37C29}
\keywords{Alexander duality, linking number, vortex lines, helicity,
capacity veto, blow-up exclusion, golden ratio}
\maketitle

%=======================================================================
\section{Introduction}\label{sec:intro}
%=======================================================================

The question addressed by this paper is topological:
\emph{Can a rigid rotation (a constant-vorticity flow with parallel vortex
lines extending over all of $\R^3$) arise as the rescaled blow-up limit
of a finite-energy divergence-free initial datum?}

Under a positive per-crossing penalty and a crossing-count hypothesis
(made explicit as assumption~\ref{A:infinite} in Theorem~\ref{thm:veto}),
we prove the answer is \emph{no}.
The argument proceeds in five steps:
\begin{enumerate}[nosep]
\item Finite-energy data has finite helicity (Lemma~\ref{lem:finite-helicity}),
  hence finite linking complexity.
\item Rigid rotation has zero linking over infinite spatial extent
  (Lemma~\ref{lem:rigid-linking}).
\item (Hypothesis) Any transition from finite-linking data to rigid rotation
  requires infinitely many link crossings.
\item Each crossing costs at least $\ln\varphi>0$
  (Definition~\ref{def:penalty}, Proposition~\ref{prop:penalty-pos}).
\item Finite budget cannot fund infinite cost: contradiction.
\end{enumerate}

\subsection*{Context}
This result is used in the companion paper~\cite{NS} on global regularity
for the 3D Navier--Stokes equations, where it excludes the rigid-rotation
endpoint of the running-max ancient element after direction constancy has
been established.  The crossing-count hypothesis (step~3) is verified in
that paper's blow-up framework.

%=======================================================================
\section{Alexander duality and linking in dimension three}\label{sec:alex}
%=======================================================================

We recall the topological fact that makes linking a purely
three-dimensional phenomenon.

\begin{theorem}[Alexander duality for circle complements]%
\label{thm:alexander}
Let $K\hookrightarrow S^D$ be a tamely embedded circle ($K\cong S^1$) in
the $D$-dimensional sphere.  Then
\[
H_1(S^D\setminus K;\Z)\;\cong\;\widetilde{H}^{D-2}(K;\Z)
\;\cong\;\widetilde{H}^{D-2}(S^1;\Z).
\]
The right-hand side is $\Z$ if $D-2=1$ (i.e.\ $D=3$), and trivial
otherwise.
\end{theorem}
\begin{proof}
This is classical Alexander duality; see Hatcher~\cite[\S 3.3,
Theorem~3.44]{Hatcher} or Bredon~\cite[\S VI.8]{Bredon}.
The cohomology $\widetilde{H}^k(S^1;\Z)$ is $\Z$ for $k=1$ and
$0$ for all other $k\ge0$, by the standard computation of the
cohomology of spheres.
\end{proof}

\begin{corollary}[Integer linking exists iff $D=3$]\label{cor:D3}
An integer-valued linking number for a pair of disjoint oriented closed
curves in $S^D$ (or equivalently $\R^D$) exists if and only if $D=3$.
\end{corollary}
\begin{proof}
The linking number $\mathrm{lk}(\gamma_1,\gamma_2)\in\Z$ is defined as
the homology class of $[\gamma_1]\in H_1(S^D\setminus\gamma_2;\Z)$.
By Theorem~\ref{thm:alexander}, this group is $\Z$ precisely when $D=3$,
giving a well-defined integer invariant.
When $D\neq 3$, the group is trivial and no non-trivial linking number
exists.
\end{proof}

\begin{remark}[Why dimension matters]
The special role of $D=3$ is not merely technical.
In $D=2$, curves are zero-dimensional complements (points cannot
``link'' a curve).  In $D\ge4$, disjoint curves can be unlinked by
small perturbations in the extra dimensions.  Only in $D=3$ is the
complement of a curve a non-trivial knot complement with first homology
$\Z$.  See Rolfsen~\cite{Rolfsen} for an extensive treatment.
\end{remark}

%=======================================================================
\section{Helicity and linking for divergence-free fields}\label{sec:helicity}
%=======================================================================

\begin{definition}[Helicity]\label{def:helicity}
For a smooth divergence-free field $u:\R^3\to\R^3$ with
$\omega:=\nabla\times u$, the \emph{helicity} is
\[
\mathcal{H}(u):=\int_{\R^3}u\cdot\omega\,dx.
\]
\end{definition}

\begin{lemma}[Finite helicity from $H^1$ data]\label{lem:finite-helicity}
If $u_0\in H^1(\R^3;\R^3)$ is divergence-free, then
\[
|\mathcal{H}(u_0)|
\le \|u_0\|_{L^2}\,\|\omega_0\|_{L^2}
\le \|u_0\|_{H^1}^2
<\infty.
\]
\end{lemma}
\begin{proof}
Set $\omega_0:=\nabla\times u_0$.  By the Cauchy--Schwarz inequality in
$L^2(\R^3)$:
\[
|\mathcal{H}(u_0)|
=\left|\int_{\R^3}u_0\cdot\omega_0\,dx\right|
\le \|u_0\|_{L^2}\,\|\omega_0\|_{L^2}.
\]
For the second inequality, note that
$\|\omega_0\|_{L^2}=\|\nabla\times u_0\|_{L^2}
\le\|\nabla u_0\|_{L^2}\le\|u_0\|_{H^1}$
(since $|\nabla\times u|\le|\nabla u|$ pointwise).
Also $\|u_0\|_{L^2}\le\|u_0\|_{H^1}$.
Hence $|\mathcal{H}(u_0)|\le\|u_0\|_{H^1}^2<\infty$.
\end{proof}

\begin{remark}[Arnol'd's helicity-linking interpretation]
The helicity $\mathcal{H}$ equals the average asymptotic linking
number of pairs of vortex lines, weighted by vorticity magnitude.
This was established by Arnol'd~\cite{Arnold} and is developed in
detail in Arnol'd--Khesin~\cite[\S III.1]{ArnoldKhesin}.
Finite helicity therefore implies a finite ``total linking complexity''
of the initial data.
\end{remark}

%=======================================================================
\section{Rigid rotation: zero linking over infinite extent}\label{sec:rigid}
%=======================================================================

\begin{definition}[Rigid rotation]\label{def:rigid}
The \emph{rigid rotation} in $\R^3$ is the velocity field
$u_{\rm rig}(x)=\frac{1}{2}(-x_2,x_1,0)$,
with constant vorticity $\omega_{\rm rig}=(0,0,1)$.
\end{definition}

\begin{lemma}[Vortex lines and linking]\label{lem:rigid-linking}
The vortex lines of the rigid rotation are the vertical lines
$\ell_{a,b}:=\{(a,b,t):t\in\R\}$ for $(a,b)\in\R^2$.
Every pair of distinct vortex lines has linking number~$0$.
\end{lemma}
\begin{proof}
The vortex lines are the integral curves of
$\omega_{\rm rig}=(0,0,1)$, which are the vertical lines parallel to
the $x_3$-axis.
Two parallel lines in $\R^3$ can be separated by a plane; by the
isotopy classification of links (see Rolfsen~\cite[\S 5A]{Rolfsen}),
parallel lines are unlinked, hence $\mathrm{lk}(\ell_{a,b},\ell_{a',b'})=0$
for all $(a,b)\neq(a',b')$.
\end{proof}

\begin{lemma}[Infinite energy of rigid rotation]\label{lem:rigid-infinite}
$\|u_{\rm rig}\|_{L^2(\R^3)}=\infty$.
\end{lemma}
\begin{proof}
$|u_{\rm rig}(x)|^2=(x_1^2+x_2^2)/4$.
Integrating over $\R^3$:
\[
\int_{\R^3}\frac{x_1^2+x_2^2}{4}\,dx
=\frac14\int_{-\infty}^\infty\!\int_{\R^2}(x_1^2+x_2^2)\,dx_1\,dx_2\,dx_3
=\infty,
\]
since $\int_{\R^2}(x_1^2+x_2^2)\,dx_1\,dx_2=\infty$ and the $x_3$
integral contributes an additional infinite factor.
\end{proof}

%=======================================================================
\section{The link penalty}\label{sec:penalty}
%=======================================================================

\begin{definition}[Link penalty]\label{def:penalty}
Each \emph{topological crossing event} (a change of linking number by
$\pm 1$ between a pair of vortex lines) incurs a cost
\[
\Delta J:=\ln\varphi,
\]
where $\varphi=(1+\sqrt{5})/2$ is the golden ratio.
\end{definition}

\begin{proposition}[Positivity of the link penalty]\label{prop:penalty-pos}
$\Delta J=\ln\varphi>0$.
\end{proposition}
\begin{proof}
$\varphi=(1+\sqrt5)/2>(1+2)/2=3/2>1$, hence $\ln\varphi>0$.
Numerically, $\ln\varphi\approx 0.4812$.
\end{proof}

\begin{remark}[Origin of the penalty]
The value $\ln\varphi$ is the minimal nonzero cost on the self-similar
scale lattice of the reciprocal cost function $J$;
see~\cite[Definition~6.2 and Proposition~6.3]{F1} for the derivation.
The golden ratio arises from the self-similarity equation
$\varphi^2=\varphi+1$, which characterizes the unique positive
fixed point of the quadratic iteration.
\end{remark}

%=======================================================================
\section{The finite-capacity veto}\label{sec:veto}
%=======================================================================

\begin{theorem}[Conditional finite-capacity veto]\label{thm:veto}
Let $u_0\in H^1(\R^3;\R^3)$ be divergence-free.  Assume:
\begin{enumerate}[label=\textup{(A\arabic*)},nosep]
\item\label{A:penalty}
  \textbf{Positive per-crossing cost:} every topological crossing event
  changes linking by at most $\pm 1$ and incurs cost at least
  $\Delta J=\ln\varphi>0$.
\item\label{A:infinite}
  \textbf{Infinite crossing count:} any admissible evolution from $u_0$
  to the rigid rotation profile $u_{\rm rig}$ requires infinitely many
  crossing events.
\item\label{A:budget}
  \textbf{Finite topological budget:} the total crossing cost available
  in the admissible evolution class is bounded by a finite quantity
  determined by $\|u_0\|_{H^1}$.
\end{enumerate}
Then no admissible evolution from $u_0$ can reach rigid rotation.
\end{theorem}
\begin{proof}
\textbf{Step 1 (Cost per crossing).}
By~\ref{A:penalty}, each crossing event costs at least
$\Delta J=\ln\varphi>0$.

\textbf{Step 2 (Infinite crossings needed).}
By~\ref{A:infinite}, reaching rigid rotation from $u_0$ requires
infinitely many crossings, say $N=\infty$.

\textbf{Step 3 (Infinite total cost).}
The total crossing cost is bounded below by
\[
\sum_{n=1}^N \Delta J
=\sum_{n=1}^\infty \ln\varphi
=\infty.
\]

\textbf{Step 4 (Budget exhaustion).}
By~\ref{A:budget}, the available budget is some finite $B<\infty$.
But Step~3 shows the required cost is $\infty>B$.

\textbf{Conclusion.}
The required cost exceeds the available budget, so no admissible
evolution reaches rigid rotation.
\end{proof}

\begin{remark}[Role of the crossing-count hypothesis]\label{rem:crossing}
Assumption~\ref{A:infinite} is the topological crossing-count input.
It is \emph{not} proved in this paper; its verification is external
and depends on the specific blow-up analysis framework.
In the companion Navier--Stokes paper~\cite{NS}, this hypothesis is
verified using the direction-constancy mechanism and the Arnol'd
helicity--linking correspondence.

The other two assumptions are unconditional:
\ref{A:penalty} follows from the positivity of $\ln\varphi$
(Proposition~\ref{prop:penalty-pos}), and
\ref{A:budget} follows from the finiteness of helicity
(Lemma~\ref{lem:finite-helicity}).
\end{remark}

\begin{corollary}[Application to Navier--Stokes blow-up exclusion]
\label{cor:ns}
In the running-max blow-up analysis of the 3D incompressible
Navier--Stokes equations, suppose direction constancy forces the ancient
element into the rigid-rotation class and
assumptions~\ref{A:penalty}--\ref{A:budget} are verified in that
framework.
Then Theorem~\ref{thm:veto} provides the final contradiction:
the ancient element must be trivial ($\omega^\infty\equiv 0$),
contradicting the running-max normalization $|\omega^\infty(0,0)|=1$.
See~\cite{NS} for the complete argument.
\end{corollary}

%=======================================================================
\section*{Acknowledgments}
%=======================================================================
The author thanks the referees for careful reading and for suggesting
the explicit separation of unconditional and conditional ingredients.

\begin{thebibliography}{99}

\bibitem{Arnold}
V.~I.~Arnol'd,
\emph{The asymptotic Hopf invariant and its applications},
Selecta Math.\ Soviet.\ \textbf{5} (1986), 327--345.
(Original Russian: 1974.)

\bibitem{ArnoldKhesin}
V.~I.~Arnol'd and B.~A.~Khesin,
\emph{Topological Methods in Hydrodynamics},
Applied Mathematical Sciences, vol.~125, Springer, 1998.

\bibitem{Bredon}
G.~E.~Bredon,
\emph{Topology and Geometry},
Graduate Texts in Mathematics, vol.~139, Springer, 1993.

\bibitem{F1}
J.~Washburn,
\emph{The canonical reciprocal cost: exact multi-bond minimization,
frustration lower bounds, and simultaneous-vs-sequential descent},
preprint, 2026.

\bibitem{Hatcher}
A.~Hatcher,
\emph{Algebraic Topology},
Cambridge Univ.\ Press, 2002.

\bibitem{NS}
J.~Washburn,
\emph{Global regularity for the 3D incompressible Navier--Stokes
equations via the topological frustration pinch},
preprint, 2026.

\bibitem{Rolfsen}
D.~Rolfsen,
\emph{Knots and Links},
Mathematics Lecture Series, vol.~7, Publish or Perish, 1976.
Corrected reprint: AMS Chelsea, 2003.

\end{thebibliography}

\end{document}
